\documentclass{article}
\usepackage{multicol}
\usepackage{graphicx}
\usepackage{subcaption}
\usepackage{eso-pic}
\usepackage{graphicx,float}
\usepackage[top=3.2cm, bottom=0cm, outer=2cm, inner=2cm]{geometry}

\title{\textbf{Uma An\'{a}lise sobre Daltonismo e Realidade Virtual}}
\author{\textbf{Jo\~{a}o Baptista Assun\c{c}\~{a}o Pereira e Silva, Rosilane Ribeiro da Mota}\\ Departamento de Teatro e Cinema. Universidade de Minas Gerais, Belo Horizonte, Brasil}
\date{}

\newcommand\BackgroundIm{\put(0,0){\parbox[b][\paperheight]{\paperwidth}{\vfill\centering         \includegraphics[height=\paperheight,width=\paperwidth]{Recursos/fundo.png}\vfill}}}

\begin{document}
    \AddToShipoutPicture*{\BackgroundIm}
    \maketitle
    \begin{multicols}{2}
            \section*{Introdu\c{c}\~{a}o}
        O avan\c{c}o recente nas \'{a}reas da realidade virtual e tecnologias
        de imers\~{a}o permite que os usu\'{a}rios de ditas tecnologias se
        coloquem em experi\^{e}ncias que expandem a realidade, que
        permitem que o potencial criativo de uma pessoa possa ser
        utilizado de formas n\~{a}o convencionais para trabalho, ou
        simples experimenta\c{c}\~{a}o. O uso dos aparatos de realidade
        virtual requer certas caracter\'{\i}sticas fision\^{o}micas que s\~{a}o uma
        necessidade para a experi\^{e}ncia de imers\~{a}o e infelizmente nem
        todos os usu\'{a}rios podem cumprir totalmente com algumas
        dessas caracter\'{\i}sticas. Para lidar com essas quest\~{o}es na \'{a}rea
        de realidade virtual, games, e v\'{a}rias formas de m\'{\i}dia visual,
        foram pesquisadas v\'{a}rias formas que podem auxiliar o
        portador de disfun\c{c}\~{o}es.
        
            \section*{Metodologia}
        An\'{a}lise de artigos, livros e criticas sobre os elementos
        trabalhados. Tendo enfoque em artigos relacionados \`{a}s
        disfun\c{c}\~{o}es visuais colocando o portador usu\'{a}rio de softwares
        que dependem da vis\~{a}o para imers\~{a}o como foco. Este
        levantamento leva em conta as dificuldades que os dalt\^{o}nicos
        sofrem para se integrar na sociedade como membros
        funcionais e como essas dificuldades poderiam se apresentar
        em meios de realidade virtual. Essa identifica\c{c}\~{a}o foi feita a
        partir de artigos que tratam da \'{a}rea de desenvolvimento de
        tecnologias para meio digital, desenvolvimento de intera\c{c}\~{a}o
        com realidade aumentada e realidade virtual e \'{a}reas voltada
        para adapta\c{c}\~{a}o de produtos e servi\c{c}os para dalt\^{o}nicos entre e
        demais portadores de disfun\c{c}\~{o}es visuais relacionadas \`{a} cor. A
        partir dessa identifica\c{c}\~{a}o foi feito um cruzamento de
        informa\c{c}\~{o}es com o intuito de levantar quest\~{o}es pertinentes
        sobre as \'{a}reas estudadas.
        
            \section*{Discuss\~{a}o}
        A presen\c{c}a e o realismo s\~{a}o dois elementos que s\~{a}o
        relacionados diretamente ao ambiente de realidade virtual. A
        presen\c{c}a ou imers\~{a}o \'{e} um dos elementos-chave para cria\c{c}\~{a}o
        da experi\^{e}ncia imersiva no ambiente de realidade virtual. A
        rela\c{c}\~{a}o entre os dois fatores \'{e} derivada do qu\~{a}o
        geometricamente realista o ambiente \'{e}, sendo que quanto
        maior o grau de realismo, maior a presen\c{c}a ou imers\~{a}o.
        games, e v\'{a}rias formas de m\'{\i}dia visual, foram pesquisadas
        v\'{a}rias formas que podem auxiliar o portador de disfun\c{c}\~{o}es.
        Alguns jogos possuem recursos que ajustam a paleta de cores
        para ajudar no contraste entre tons semelhantes, permitindo
        que seus usu\'{a}rios dalt\^{o}nicos, ou cegos para cor possam
        distinguir facilmente as informa\c{c}\~{o}es que recebem, facilitando
        a compreens\~{a}o.
        
       \begin{figure}[H]
            \centering
            \includegraphics[width=4cm]{Recursos/oculos.png}
            \caption*{Figura 1. Occulus Rift (esquerda) e HTC Vive (Direita)}
        \end{figure}
         
            \section*{Conclus\~{a}o}
        A popula\c{c}\~{a}o afetada por defici\^{e}ncias relacionadas \`{a} vis\~{a}o que
        de alguma forma alteram a experi\^{e}ncia visual possui recursos
        limitados para lidar com essas defici\^{e}ncias tanto dentro,
        quanto fora do meio digital de maneira que esses recursos
        podem ser mais facilmente ajustados para o usu\'{a}rio. Esses
        m\'{e}todos de ajuste n\~{a}o variam muito e tem uma base
        semelhante tanto na realidade virtual quanto no meio f\'{\i}sic,
        assim seus desenvolvedores chegaram a meios semelhantes
        para atenuar o problema tanto no mundo real, quanto no
        ambiente virtual, e por consequ\^{e}ncia seguem a mesma base
        de requisitos para que o m\'{e}todo de ajuste funcione.
    \end{multicols}
\end{document}
